\chapter{Auxiliary Theorems and Proofs}
\label{ch:appendix}

\providecommand{\Q}{\mathbb{Q}}
\providecommand{\GL}{\operatorname{GL}}
\providecommand{\SO}{\operatorname{SO}}
\providecommand{\Orth}{\operatorname{O}}
\providecommand{\SU}{\operatorname{SU}}
\providecommand{\Sym}{\operatorname{Sym}}
\providecommand{\rank}{\operatorname{rank}}
\providecommand{\tr}{\operatorname{tr}}
\providecommand{\Img}{\operatorname{Im}}
\providecommand{\Ric}{\operatorname{Ric}}
\providecommand{\Vol}{\operatorname{Vol}}
\providecommand{\diag}{\operatorname{diag}}
\providecommand{\II}{\mathrm{II}}
\providecommand{\dvol}{\,dA}

This appendix collects classical theorems of differential geometry, linear
algebra, group theory, measure theory, and functional analysis, together with
their proofs. Although these results are fundamental to the mathematical rigour
of the main text, they are not central to the theory of geometric deep learning
developed in the principal chapters, and gathering them here keeps the exposition
focused. Each section presents a theorem with the preliminaries it requires, a
formal statement, and a complete proof or a proof sketch with references to the
original sources. The measure-theoretic and functional-analytic machinery in the
final two sections supports the universal approximation theorem, whose statement
remains in \Cref{ch:foundations} and which refers back to this appendix for its
heavy prerequisites.

%=======================================================================
\section{The Gauss--Bonnet Theorem}
\label{sec:app:gauss-bonnet}
%=======================================================================

The Gauss--Bonnet theorem is one of the deepest results in differential
geometry, establishing a fundamental link between the local geometry (curvature)
and the global topology (Euler characteristic) of a surface. It shows that
certain geometric quantities are determined solely by the topology of the
surface.

\subsection{Preliminaries}
\label{subsec:app:gauss-bonnet-prelim}

Before stating the theorem we introduce the necessary notions.

\begin{definition}[Gaussian curvature]
\label{def:app:gaussian-curvature}
Let $\mathcal{M}$ be a regular surface immersed in $\R^3$. At each point
$p \in \mathcal{M}$ consider the \emph{principal curvatures} $\kappa_1$ and
$\kappa_2$, the eigenvalues of the second fundamental form. The \emph{Gaussian
curvature} $K \colon \mathcal{M} \to \R$ is the product of the principal
curvatures:
\begin{equation}
\label{eq:app:gaussian-curvature}
K = \kappa_1 \kappa_2 .
\end{equation}
\end{definition}

\begin{remark}
Gauss's \emph{Theorema Egregium} states that the Gaussian curvature $K$ is an
\emph{isometric invariant}: it depends only on the first fundamental form (the
metric) and not on how the surface is immersed in $\R^3$. Thus $K$ is an
intrinsic property of the surface.
\end{remark}

\begin{definition}[Geodesic curvature]
\label{def:app:geodesic-curvature}
Let $\gamma \colon [a,b] \to \mathcal{M}$ be a regular curve parametrized by arc
length on an oriented surface $\mathcal{M}$. Let $\mathbf{n}$ be the unit normal
to the surface and $\mathbf{T} = \gamma'$ the unit tangent to the curve. The
\emph{geodesic curvature} $\kappa_g$ of $\gamma$ is
\begin{equation}
\label{eq:app:geodesic-curvature}
\kappa_g = \inner{\gamma''}{\mathbf{n} \times \mathbf{T}} ,
\end{equation}
where $\mathbf{n} \times \mathbf{T}$ is the normal to the curve lying in the
tangent plane to the surface.
\end{definition}

\begin{remark}
The geodesic curvature $\kappa_g$ measures how far a curve deviates from being a
geodesic within the surface. A curve is a geodesic if and only if $\kappa_g = 0$
at every point. Intuitively, $\kappa_g$ captures the component of curvature that
is ``visible'' from within the surface.
\end{remark}

\begin{definition}[Euler characteristic]
\label{def:app:euler-characteristic}
Let $\mathcal{M}$ be a compact surface. The \emph{Euler characteristic}
$\chi(\mathcal{M})$ is a topological invariant defined equivalently as follows:
\begin{enumerate}
    \item \textbf{Triangulation.} If $\mathcal{M}$ admits a triangulation with
          $V$ vertices, $E$ edges, and $F$ faces, then
          \begin{equation}
          \label{eq:app:euler-triangulation}
          \chi(\mathcal{M}) = V - E + F .
          \end{equation}
    \item \textbf{Genus.} For an orientable surface of genus $g$ (number of
          ``handles''),
          \begin{equation}
          \label{eq:app:euler-genus}
          \chi(\mathcal{M}) = 2 - 2g .
          \end{equation}
\end{enumerate}
\end{definition}

\begin{example}
\label{ex:app:euler-surfaces}
Some values of the Euler characteristic for familiar surfaces:
\begin{itemize}
    \item Sphere $S^2$: $\chi(S^2) = 2$ (genus $g = 0$).
    \item Torus $T^2$: $\chi(T^2) = 0$ (genus $g = 1$).
    \item Double torus: $\chi = -2$ (genus $g = 2$).
    \item Disc: $\chi = 1$ (a surface with boundary).
\end{itemize}
\end{example}

\subsection{Statement of the theorem}
\label{subsec:app:gauss-bonnet-statement}

We present first the general version for surfaces with boundary and then derive
important special cases.

\begin{theorem}[Gauss--Bonnet for surfaces with boundary]
\label{thm:app:gauss-bonnet-general}
Let $\mathcal{M}$ be a compact oriented Riemannian surface with piecewise-smooth
boundary $\partial \mathcal{M}$. Suppose $\partial \mathcal{M}$ consists of
regular curves meeting at vertices where they form exterior angles
$\alpha_1, \ldots, \alpha_n$. Then
\begin{equation}
\label{eq:app:gauss-bonnet-general}
\int_{\mathcal{M}} K \, dA + \int_{\partial \mathcal{M}} \kappa_g \, ds
+ \sum_{i=1}^n \alpha_i = 2\pi \chi(\mathcal{M}) ,
\end{equation}
where $K$ is the Gaussian curvature, $dA$ the area element induced by the
Riemannian metric, $\kappa_g$ the geodesic curvature of the boundary, $ds$ the
arc-length element, and $\alpha_i$ the exterior angle at the $i$-th boundary
vertex.
\end{theorem}

\begin{corollary}[Gauss--Bonnet for closed surfaces]
\label{cor:app:gauss-bonnet-closed}
If $\mathcal{M}$ is a compact surface without boundary (closed), then
\begin{equation}
\label{eq:app:gauss-bonnet-closed}
\int_{\mathcal{M}} K \, dA = 2\pi \chi(\mathcal{M}) .
\end{equation}
\end{corollary}

\begin{proof}
If $\partial \mathcal{M} = \emptyset$, the boundary integral and the sum of
exterior angles both vanish, and the result follows directly from
\Cref{thm:app:gauss-bonnet-general}.
\end{proof}

\begin{corollary}[Angle sum in geodesic triangles]
\label{cor:app:gauss-bonnet-triangle}
Let $T \subset \mathcal{M}$ be a geodesic triangle (a region bounded by three
geodesics) on a Riemannian surface, with interior angles $\alpha$, $\beta$,
$\gamma$. Then
\begin{equation}
\label{eq:app:gauss-bonnet-triangle}
\alpha + \beta + \gamma = \pi + \int_T K \, dA .
\end{equation}
\end{corollary}

\begin{proof}
Apply \Cref{thm:app:gauss-bonnet-general} to the triangle $T$:
\begin{enumerate}
    \item Since $T$ is homeomorphic to a disc, $\chi(T) = 1$.
    \item Since the sides are geodesics, $\kappa_g = 0$ along the whole
          boundary, so $\int_{\partial T} \kappa_g \, ds = 0$.
    \item The exterior angles at the vertices satisfy
          $\alpha_i^{\text{ext}} = \pi - \alpha_i^{\text{int}}$, where
          $\alpha_i^{\text{int}}$ is the corresponding interior angle.
\end{enumerate}
The theorem therefore gives
\[
\int_T K \, dA + 0 + \sum_{i=1}^3 (\pi - \theta_i) = 2\pi ,
\]
where $\theta_1, \theta_2, \theta_3$ are the interior angles
$\alpha, \beta, \gamma$. Simplifying,
\[
\int_T K \, dA + 3\pi - (\alpha + \beta + \gamma) = 2\pi ,
\qquad\text{hence}\qquad
\alpha + \beta + \gamma = \pi + \int_T K \, dA . \qedhere
\]
\end{proof}

\begin{remark}[Geometric interpretation]
\Cref{cor:app:gauss-bonnet-triangle} has profound implications:
\begin{itemize}
    \item \textbf{Euclidean geometry} ($K = 0$): $\alpha + \beta + \gamma = \pi$
          (the classical result of Euclid).
    \item \textbf{Spherical geometry} ($K > 0$):
          $\alpha + \beta + \gamma > \pi$ (angular excess proportional to area).
    \item \textbf{Hyperbolic geometry} ($K < 0$):
          $\alpha + \beta + \gamma < \pi$ (angular defect proportional to area).
\end{itemize}
This shows that Euclid's fifth (parallel) postulate is equivalent to the
condition $K = 0$.
\end{remark}

\subsection{Proof of the theorem}
\label{subsec:app:gauss-bonnet-proof}

We present a proof sketch following the classical
approach~\citep{do1976differential,lee2018introduction}.

\begin{lemma}[Formula for geodesic polygons]
\label{lem:app:geodesic-polygon}
Let $P$ be a simple geodesic polygon on a surface $\mathcal{M}$ with $n$ vertices
and interior angles $\theta_1, \ldots, \theta_n$. Then
\begin{equation}
\label{eq:app:geodesic-polygon}
\sum_{i=1}^n \theta_i = (n-2)\pi + \int_P K \, dA .
\end{equation}
\end{lemma}

\begin{proof}
Triangulate the polygon $P$ into $n-2$ geodesic triangles
$T_1, \ldots, T_{n-2}$ (always possible for simple polygons). Applying
\Cref{cor:app:gauss-bonnet-triangle} to each triangle and summing,
\[
\sum_{j=1}^{n-2} (\alpha_j + \beta_j + \gamma_j)
= (n-2)\pi + \sum_{j=1}^{n-2} \int_{T_j} K \, dA .
\]
On the left-hand side, the angles at interior vertices each sum to $2\pi$ (a full
angle), while the angles at boundary vertices are exactly the $\theta_i$ of the
original polygon. On the right-hand side, the sum of integrals is the integral
over all of $P$:
\[
\sum_{i=1}^n \theta_i + 2\pi \cdot (\text{interior vertices})
= (n-2)\pi + \int_P K \, dA .
\]
A careful count of the triangulation (using Euler's formula for planar graphs)
shows that the contribution of the interior vertices cancels, yielding the
desired result.
\end{proof}

\begin{proof}[Proof sketch of \Cref{thm:app:gauss-bonnet-general}]
The proof proceeds in several steps.

\textbf{Step 1: Triangulation of the surface.}
By \Cref{thm:app:rado}, every compact surface admits a triangulation. Refining
if necessary, we obtain a geodesic triangulation
$\mathcal{T} = \{T_1, \ldots, T_F\}$ in which every triangle has geodesic sides.

\textbf{Step 2: Apply the lemma to each triangle.}
Apply \Cref{lem:app:geodesic-polygon} (case $n = 3$) to each triangle:
\[
\alpha_j + \beta_j + \gamma_j = \pi + \int_{T_j} K \, dA ,
\qquad j = 1, \ldots, F .
\]

\textbf{Step 3: Sum over all triangles.}
Summing over all triangles,
\[
\sum_{j=1}^F (\alpha_j + \beta_j + \gamma_j) = F\pi + \int_{\mathcal{M}} K \, dA .
\]

\textbf{Step 4: Analysis of the angle sum.}
The angle sum on the left decomposes according to the location of the vertices:
\begin{itemize}
    \item \textit{Interior vertices}: each contributes $2\pi$ (a full angle
          around the point). If there are $V_{\text{int}}$ interior vertices,
          they contribute $2\pi V_{\text{int}}$.
    \item \textit{Smooth boundary vertices}: each on a smooth part of the
          boundary contributes $\pi$ (a half angle).
    \item \textit{Corner boundary vertices}: each corner with interior angle
          $\theta_i$ contributes exactly $\theta_i = \pi - \alpha_i$.
\end{itemize}

\textbf{Step 5: Apply Euler's formula.}
Let $V$, $E$, $F$ denote the total numbers of vertices, edges, and faces in the
triangulation. By Euler's formula, $V - E + F = \chi(\mathcal{M})$. Moreover,
counting edges: each triangle has three edges, each interior edge is shared by
two triangles, and boundary edges belong to a single triangle, which yields
relations among $V$, $E$, $F$ and the boundary quantities.

\textbf{Step 6: Combine terms.}
Combining the identities of Steps 3--5 with careful algebra, the terms rearrange
to give
\[
\int_{\mathcal{M}} K \, dA + \int_{\partial \mathcal{M}} \kappa_g \, ds
+ \sum_{i=1}^n \alpha_i = 2\pi \chi(\mathcal{M}) .
\]
The geodesic-curvature integral appears naturally on passing to the limit as the
triangulation is refined near the boundary.
\end{proof}

\begin{remark}
A complete and rigorous proof can be found in \citet[Chapter~4]{do1976differential}
or in \citet[Chapter~9]{lee2018introduction}. The version for abstract Riemannian
manifolds (not necessarily immersed in $\R^3$) requires connections and parallel
transport.
\end{remark}

\subsection{Generalizations}
\label{subsec:app:gauss-bonnet-generalizations}

The Gauss--Bonnet theorem admits several important generalizations:
\begin{enumerate}
    \item \textbf{Gauss--Bonnet--Chern theorem.} For Riemannian manifolds of
          even dimension $2n$, the theorem generalizes using the Pfaffian form of
          the curvature~\citep{chern1944simple}.
    \item \textbf{Atiyah--Singer index theorem.} A deep generalization
          connecting analytic invariants (the index of elliptic operators) with
          topological invariants.
    \item \textbf{Discrete version.} For polygonal meshes there is a discrete
          version of the theorem, useful in computational geometry and geometry
          processing.
\end{enumerate}

%=======================================================================
\section{Sard's Theorem}
\label{sec:app:sard}
%=======================================================================

Sard's theorem is a fundamental tool of analysis on manifolds that enables
``general position'' arguments: it guarantees that, in a precise sense, most
points in the codomain of a smooth map are regular values. This result is
essential in the proof of the Whitney embedding theorem.

\begin{definition}[Critical point and critical value]
\label{def:app:critical-point}
Let $F \colon \mathcal{M} \to \mathcal{N}$ be a smooth map between smooth
manifolds. A point $p \in \mathcal{M}$ is a \emph{critical point} of $F$ if the
differential $dF_p \colon T_p\mathcal{M} \to T_{F(p)}\mathcal{N}$ is not
surjective. A point $y \in \mathcal{N}$ is a \emph{critical value} if it is the
image of some critical point, that is, if $y = F(p)$ for some critical point $p$.
\end{definition}

\begin{definition}[Regular value]
\label{def:app:regular-value}
A point $y \in \mathcal{N}$ is a \emph{regular value} of
$F \colon \mathcal{M} \to \mathcal{N}$ if it is not a critical value, that is, if
for every $p \in F^{-1}(y)$ the differential $dF_p$ is surjective. In particular,
if $y \notin F(\mathcal{M})$, then $y$ is trivially a regular value.
\end{definition}

\begin{theorem}[Sard's theorem]
\label{thm:app:sard}
Let $F \colon \mathcal{M} \to \mathcal{N}$ be a smooth map between smooth
manifolds. Then the set of critical values of $F$ has Lebesgue measure zero in
$\mathcal{N}$. Equivalently, the set of regular values is dense in $\mathcal{N}$.
\end{theorem}

The complete proof, which proceeds by induction on the dimension and uses measure
estimates based on the Taylor expansion of $F$, can be found in
\citet[Chapter~6, Theorem~6.10]{lee2012introduction} or in
\citet[Chapter~3]{milnor1965topology}.

\begin{remark}
The power of Sard's theorem is that it lets one conclude the existence of objects
with good properties (regular values) without constructing them explicitly: it
suffices to show that the ``bad'' ones form a set of measure zero. In the proof
of Whitney's theorem this argument guarantees the existence of projection
directions that preserve the embedding.
\end{remark}

%=======================================================================
\section{The Regular Level Set Theorem}
\label{sec:app:regular-level}
%=======================================================================

The regular level set theorem connects the theory of smooth maps between
manifolds with the submanifold structure of their level sets. It is a
fundamental tool of differential geometry: it constructs submanifolds as
preimages of regular values, and it is the theoretical basis for the structure of
the fibres of a submersion.

\begin{theorem}[Regular level set theorem]
\label{thm:app:regular-level}
Let $F \colon \mathcal{M} \to \mathcal{N}$ be a smooth map between smooth
manifolds of dimensions $m = \dim \mathcal{M}$ and $n = \dim \mathcal{N}$, with
$m \geq n$. If $q \in \mathcal{N}$ is a regular value of $F$
(\Cref{def:app:regular-value}) with $F^{-1}(q) \neq \emptyset$, then $F^{-1}(q)$
is a smooth submanifold of $\mathcal{M}$ of dimension $m - n$.
\end{theorem}

\begin{proof}
Let $p \in F^{-1}(q)$. Since $q$ is a regular value, the differential
$dF_p \colon T_p\mathcal{M} \to T_q\mathcal{N}$ is surjective, hence of rank $n$.
Choose charts $(U, \varphi)$ centred at $p$ and $(V, \psi)$ centred at $q$ with
$F(U) \subseteq V$. In these coordinates the local representation
$\hat{F} = \psi \circ F \circ \varphi^{-1} \colon \varphi(U) \subseteq \R^m \to
\R^n$ has Jacobian $D\hat{F}_{\varphi(p)}$ of rank $n$.

Reordering coordinates if necessary, we may assume the last $n$ columns of
$D\hat{F}_{\varphi(p)}$ form an invertible $n \times n$ block. Write
$\R^m = \R^{m-n} \times \R^n$ with coordinates $(x', x'')$, where
$x' = (x_1, \ldots, x_{m-n})$ and $x'' = (x_{m-n+1}, \ldots, x_m)$.

Consider the map $\Phi \colon \varphi(U) \to \R^m$ given by
$\Phi(x', x'') = (x', \hat{F}(x', x''))$. Its Jacobian at $\varphi(p)$ is
\[
D\Phi_{\varphi(p)} = \begin{pmatrix} I_{m-n} & 0 \\[2pt]
\frac{\partial \hat{F}}{\partial x'} & \frac{\partial \hat{F}}{\partial x''}
\end{pmatrix} ,
\]
which is invertible, being block lower triangular with invertible diagonal blocks
$I_{m-n}$ and $\frac{\partial \hat{F}}{\partial x''}$. By the inverse function
theorem (\Cref{thm:app:inverse-function}), $\Phi$ is a local diffeomorphism on a
neighbourhood $W$ of $\varphi(p)$.

In the new coordinates $\Phi$, the map $F$ reads
\[
\hat{F} \circ \Phi^{-1}(y', y'') = y'' ,
\]
that is, $\hat{F}$ becomes the projection onto the last $n$ coordinates.
Therefore
\[
(\hat{F} \circ \Phi^{-1})^{-1}(\psi(q))
= \{(y', y'') \in \Phi(W) : y'' = \psi(q)\}
= \{(y', \psi(q)) : y' \in U'\}
\]
for some open $U' \subseteq \R^{m-n}$. This shows that, near each point $p$,
$F^{-1}(q)$ is locally the graph of a smooth function of $m - n$ variables, which
endows it with the structure of a smooth submanifold of dimension $m - n$.
\end{proof}

\begin{corollary}[Fibres of a submersion]
\label{cor:app:submersion-fibres}
If $F \colon \mathcal{M} \to \mathcal{N}$ is a smooth submersion, then for every
$q \in F(\mathcal{M})$ the fibre $F^{-1}(q)$ is a smooth submanifold of
$\mathcal{M}$ of dimension $\dim \mathcal{M} - \dim \mathcal{N}$.
\end{corollary}

\begin{proof}
Since $F$ is a submersion, $dF_p$ is surjective for every $p \in \mathcal{M}$. In
particular, every point $q \in F(\mathcal{M})$ is a regular value of $F$, and the
result follows from \Cref{thm:app:regular-level}.
\end{proof}

\begin{remark}
A complete reference for the regular level set theorem and its applications can
be found in \citet[Chapter~5, Theorem~5.12]{lee2012introduction}. The classical
formulation in $\R^n$ is obtained as a special case by taking
$\mathcal{M} = \R^m$ and $\mathcal{N} = \R^n$ with the identity charts.
\end{remark}

%=======================================================================
\section{Area of Geodesic Discs in Surfaces of Constant Curvature}
\label{sec:app:geodesic-disc}
%=======================================================================

The area of a geodesic disc of radius $R$ in a surface of constant Gaussian
curvature $K$ satisfies a formula that generalizes the Euclidean result
$\pi R^2$. We derive this formula from the Jacobi equation in geodesic polar
coordinates.

\begin{proposition}[Metric in geodesic polar coordinates]
\label{prop:app:geodesic-polar-metric}
Let $(\mathcal{M}, g)$ be a Riemannian surface of constant Gaussian curvature $K$,
and let $p \in \mathcal{M}$. In geodesic polar coordinates $(r, \theta)$ centred
at $p$ (where $r = d_g(p, \cdot)$ and $\theta \in [0, 2\pi)$ is the angle in
$T_p\mathcal{M}$), the Riemannian metric takes the form
\begin{equation}
\label{eq:app:geodesic-polar-metric}
ds^2 = dr^2 + f_K(r)^2 \, d\theta^2 ,
\end{equation}
where the function $f_K \colon [0, \infty) \to \R$ satisfies the \textbf{Jacobi
equation}
\begin{equation}
\label{eq:app:jacobi-radial}
f_K''(r) + K \, f_K(r) = 0 , \qquad f_K(0) = 0 , \quad f_K'(0) = 1 .
\end{equation}
\end{proposition}

\begin{proof}
In geodesic polar coordinates, the radial geodesics from $p$ correspond to
$\theta = \text{const}$, with natural arc-length parametrization $r$. The metric
coefficient $g_{rr} = 1$ reflects this parametrization. The angular coefficient
$g_{\theta\theta} = f_K(r)^2$ measures how adjacent radial geodesics spread
apart; the function $f_K(r)$ is the norm of the Jacobi field $J(r)$ along a
radial geodesic with initial conditions $J(0) = 0$ and $\norm{J'(0)} = 1$. The
Jacobi equation $J'' + R(\dot{\gamma}, J)\dot{\gamma} = 0$, in the constant
curvature case $K$ in dimension $2$, reduces to the scalar equation
$f_K'' + K f_K = 0$ \citep[Chapter~10]{lee2018introduction}.
\end{proof}

The solution of equation~\eqref{eq:app:jacobi-radial} depends on the sign of $K$:
\begin{equation}
\label{eq:app:fK-solutions}
f_K(r) = \begin{cases}
\dfrac{1}{\sqrt{K}} \sin(\sqrt{K} \, r) & \text{if } K > 0, \\[8pt]
r & \text{if } K = 0, \\[8pt]
\dfrac{1}{\sqrt{\abs{K}}} \sinh(\sqrt{\abs{K}} \, r) & \text{if } K < 0.
\end{cases}
\end{equation}
The verification is immediate: for $K > 0$,
$f_K(r) = \frac{1}{\sqrt{K}} \sin(\sqrt{K} r)$ satisfies $f_K(0) = 0$,
$f_K'(0) = \cos(0) = 1$, and $f_K''(r) = -K f_K(r)$. The cases $K = 0$ and
$K < 0$ are checked analogously.

\begin{theorem}[Area of a geodesic disc]
\label{thm:app:geodesic-disc-area}
Let $(\mathcal{M}, g)$ be a surface of constant Gaussian curvature $K$, and let
$B_R(p) = \{q \in \mathcal{M} : d_g(p,q) \leq R\}$ be the geodesic disc of radius
$R$ centred at $p$ (with $R$ smaller than the injectivity radius). Then
\begin{equation}
\label{eq:app:geodesic-disc-area}
A(R) = \begin{cases}
\dfrac{2\pi}{K}\bigl(1 - \cos(\sqrt{K} \, R)\bigr) & \text{if } K > 0, \\[8pt]
\pi R^2 & \text{if } K = 0, \\[8pt]
\dfrac{2\pi}{\abs{K}}\bigl(\cosh(\sqrt{\abs{K}} \, R) - 1\bigr) & \text{if } K < 0.
\end{cases}
\end{equation}
Moreover, $\lim_{K \to 0} A(R) = \pi R^2$ in all cases.
\end{theorem}

\begin{proof}
The area element in geodesic polar coordinates is
$dA = f_K(r) \, dr \, d\theta$ (by~\eqref{eq:app:geodesic-polar-metric}, since
$\sqrt{\det(g)} = f_K(r)$). Therefore
\[
A(R) = \int_0^{2\pi} \int_0^R f_K(r) \, dr \, d\theta
= 2\pi \int_0^R f_K(r) \, dr .
\]

\textbf{Case $K > 0$:}
\[
A(R) = 2\pi \int_0^R \frac{1}{\sqrt{K}} \sin(\sqrt{K} \, r) \, dr
= 2\pi \left[ \frac{-\cos(\sqrt{K} \, r)}{K} \right]_0^R
= \frac{2\pi}{K} \bigl(1 - \cos(\sqrt{K} \, R)\bigr) .
\]

\textbf{Case $K = 0$:}
\[
A(R) = 2\pi \int_0^R r \, dr = \pi R^2 .
\]

\textbf{Case $K < 0$:} writing $K = -\abs{K}$,
\[
A(R) = 2\pi \int_0^R \frac{1}{\sqrt{\abs{K}}} \sinh(\sqrt{\abs{K}} \, r) \, dr
= 2\pi \left[ \frac{\cosh(\sqrt{\abs{K}} \, r)}{\abs{K}} \right]_0^R
= \frac{2\pi}{\abs{K}} \bigl(\cosh(\sqrt{\abs{K}} \, R) - 1\bigr) .
\]

\textbf{Limit $K \to 0$:} using the Taylor expansion
$\cos(x) = 1 - \frac{x^2}{2} + \frac{x^4}{24} - \cdots$,
\[
\frac{1 - \cos(\sqrt{K} R)}{K}
= \frac{1 - \bigl(1 - \frac{KR^2}{2} + \frac{K^2 R^4}{24} - \cdots\bigr)}{K}
= \frac{R^2}{2} - \frac{K R^4}{24} + \cdots
\xrightarrow{K \to 0} \frac{R^2}{2} .
\]
Multiplying by $2\pi$ gives $\pi R^2$. The case $K < 0$ is analogous using
$\cosh(x) = 1 + \frac{x^2}{2} + \cdots$.
\end{proof}

\begin{remark}[Geometric interpretation]
\label{rem:app:geodesic-disc-interpretation}
\Cref{thm:app:geodesic-disc-area} quantifies the effect of curvature on area
growth:
\begin{itemize}
    \item \textbf{Positive curvature} ($K > 0$): $A(R) < \pi R^2$. Area grows
          more slowly than in the Euclidean case. On the sphere $S^2$ ($K = 1$),
          $A(R) = 2\pi(1 - \cos R)$; for $R = \pi/2$ (a quarter of the sphere),
          $A = 2\pi$.
    \item \textbf{Negative curvature} ($K < 0$): $A(R) > \pi R^2$. Since
          $\cosh(x) \sim \frac{1}{2}e^x$ for large $x$, area grows
          \emph{exponentially}: $A(R) \sim \frac{\pi}{\abs{K}} e^{\sqrt{\abs{K}} R}$.
    \item \textbf{Connection with the Jacobi equation.} The function $f_K(r)$
          appearing in the metric is exactly the one describing the deviation of
          parallel geodesics, showing that area growth and the
          divergence/convergence of geodesics are manifestations of the same
          geometric phenomenon.
\end{itemize}
An alternative proof can be found in
\citet[Chapter~4, Section~4.6]{do1976differential} for surfaces and in
\citet[Theorem~11.9]{lee2018introduction} for the general Riemannian case.
\end{remark}

%=======================================================================
\section{Sample Complexity on Manifolds}
\label{sec:app:sample-complexity}
%=======================================================================

The number of samples required to approximate a Lipschitz function on a manifold
scales with the intrinsic dimension $m$, not with the ambient dimension $d$. We
present the necessary preliminaries and a proof sketch based on the results of
\citet{niyogi2008finding}.

\subsection{Preliminaries}
\label{subsec:app:sample-complexity-prelim}

\begin{definition}[Covering number]
\label{def:app:covering-number}
Let $(X, d)$ be a metric space and $\epsilon > 0$. An \emph{$\epsilon$-cover}
(or $\epsilon$-net) of $X$ is a subset $S \subseteq X$ such that for every
$x \in X$ there exists $s \in S$ with $d(x, s) \leq \epsilon$. The
\emph{covering number} $\mathcal{N}(X, d, \epsilon)$ is the smallest cardinality
of an $\epsilon$-cover of $X$.
\end{definition}

\begin{remark}
\label{rem:app:covering-vs-topological}
The notion of a metric $\epsilon$-cover should not be confused with a topological
cover. An $\epsilon$-cover is a set of \emph{points} whose balls of radius
$\epsilon$ cover the space, whereas a topological cover is a family of open sets
whose union contains the space. The two are related: if $S$ is an
$\epsilon$-cover of $(X,d)$, then $\{B(s, \epsilon)\}_{s \in S}$ is an open cover
of $X$ in the topological sense.
\end{remark}

\begin{lemma}[Volumetric bound on the covering number]
\label{lem:app:volumetric-covering}
Let $(\mathcal{M}, g)$ be a compact Riemannian manifold and $\epsilon > 0$. Then
\[
\mathcal{N}(\mathcal{M}, d_g, \epsilon) \geq
\frac{\Vol(\mathcal{M})}{\sup_{p \in \mathcal{M}} \Vol(B_g(p, \epsilon))} ,
\]
where $B_g(p, \epsilon)$ denotes the geodesic ball of centre $p$ and radius
$\epsilon$.
\end{lemma}

\begin{proof}
Let $S = \{s_1, \ldots, s_N\}$ be an optimal $\epsilon$-cover, that is,
$\abs{S} = \mathcal{N}(\mathcal{M}, d_g, \epsilon)$. By definition,
$\mathcal{M} \subseteq \bigcup_{i=1}^N B_g(s_i, \epsilon)$, so that
\[
\Vol(\mathcal{M}) \leq \sum_{i=1}^N \Vol(B_g(s_i, \epsilon))
\leq N \cdot \sup_{p \in \mathcal{M}} \Vol(B_g(p, \epsilon)) .
\]
Solving for $N$ gives the claimed bound.
\end{proof}

\begin{theorem}[Bishop--G\"unther comparison]
\label{thm:app:bishop-gunther}
Let $(\mathcal{M}^m, g)$ be a complete Riemannian manifold of dimension $m$ whose
Ricci curvature satisfies $\Ric \geq (m-1)\kappa_0 \, g$ for some
$\kappa_0 \in \R$. Let $\mathbb{M}^m_{\kappa_0}$ be the model space of constant
sectional curvature $\kappa_0$. Then, for every $p \in \mathcal{M}$ and $r > 0$
(with $r < \pi/\sqrt{\kappa_0}$ if $\kappa_0 > 0$),
\[
\Vol(B_g(p, r)) \leq \Vol(B_{\kappa_0}(r)) ,
\]
where $B_{\kappa_0}(r)$ denotes the geodesic ball of radius $r$ in
$\mathbb{M}^m_{\kappa_0}$.
\end{theorem}

\begin{remark}
\label{rem:app:bishop-gunther-context}
This theorem generalizes the area comparison of
\Cref{thm:app:geodesic-disc-area}, which treats the two-dimensional case, to
arbitrary dimension. In the model space, the volume of the ball of radius $r$ is
$\Vol(B_{\kappa_0}(r)) = \omega_m \int_0^r \operatorname{sn}_{\kappa_0}(t)^{m-1}
\, dt$, where $\omega_m$ is the volume of the unit ball in $\R^m$ and
$\operatorname{sn}_{\kappa_0}$ is the comparison function
($\sin(\sqrt{\kappa_0}\,t)/\sqrt{\kappa_0}$ if $\kappa_0 > 0$, $t$ if
$\kappa_0 = 0$, $\sinh(\sqrt{\abs{\kappa_0}}\,t)/\sqrt{\abs{\kappa_0}}$ if
$\kappa_0 < 0$). For small radii,
$\Vol(B_{\kappa_0}(r)) = \omega_m r^m (1 + O(\kappa_0 r^2))$. For a complete
proof, see \citet[Theorem~11.10]{lee2018introduction}.
\end{remark}

\subsection{Proof sketch}
\label{subsec:app:sample-complexity-proof}

We sketch the proof of the manifold sample-complexity bound in three steps.

\paragraph{Step 1: Reduction to covering numbers.}
Let $f \colon \mathcal{M} \to \R$ be an $L$-Lipschitz function and let
$\{x_1, \ldots, x_N\} \subset \mathcal{M}$ be a set of samples. If we wish to
approximate $f$ uniformly with error $\epsilon$ using only the values
$f(x_1), \ldots, f(x_N)$, then for any point $q \in \mathcal{M}$ there must exist
a sample $x_i$ with $d_g(q, x_i) \leq \epsilon/L$. Indeed, if
$d_g(q, x_i) > \epsilon/L$ for all $i$, the Lipschitz condition only guarantees
$\abs{f(q) - f(x_i)} \leq L \cdot d_g(q, x_i)$, which cannot be bounded by
$\epsilon$. Hence the samples must form an $(\epsilon/L)$-cover of $\mathcal{M}$,
and at least $\mathcal{N}(\mathcal{M}, d_g, \epsilon/L)$ samples are required.

\paragraph{Step 2: Lower bound on the covering number.}
Applying \Cref{lem:app:volumetric-covering} with radius $\epsilon/L$,
\[
N \geq \mathcal{N}\!\left(\mathcal{M}, d_g, \frac{\epsilon}{L}\right)
\geq \frac{\Vol(\mathcal{M})}{\sup_{p} \Vol\!\left(B_g\!\left(p,
\frac{\epsilon}{L}\right)\right)} .
\]

\paragraph{Step 3: Volume estimate.}
The sectional curvature of $\mathcal{M}$ is bounded in absolute value by
$\kappa$, which implies $\Ric \geq -(m-1)\kappa \, g$. By the Bishop--G\"unther
theorem (\Cref{thm:app:bishop-gunther}) with $\kappa_0 = -\kappa$,
\[
\Vol\!\left(B_g\!\left(p, \frac{\epsilon}{L}\right)\right)
\leq \omega_m \left(\frac{\epsilon}{L}\right)^m
\left(1 + O\!\left(\kappa \cdot \frac{\epsilon^2}{L^2}\right)\right) .
\]
The total volume $\Vol(\mathcal{M})$ admits a lower bound depending on the reach
$\tau$ and on $\kappa$ (by results of \citet{niyogi2008finding} on the geometry
of the reach). Combining,
\[
N \geq \frac{\Vol(\mathcal{M})}{\omega_m (\epsilon/L)^m
(1 + O(\kappa \epsilon^2/L^2))}
= \Omega\!\left(\left(\frac{L}{\epsilon}\right)^m \cdot
\operatorname{poly}(\tau^{-1}, \kappa)\right) .
\]
Absorbing the Lipschitz constant $L$ into the $\operatorname{poly}$ factor yields
the bound.

\begin{remark}[Intrinsic vs.\ ambient dimension]
\label{rem:app:intrinsic-complexity}
The key result is that the exponent in $1/\epsilon$ is the intrinsic dimension
$m$, not the ambient dimension $d$. For data in $\R^d$ that lie on an
$m$-dimensional manifold with $m \ll d$, the bound $\Omega((1/\epsilon)^m)$ is
exponentially smaller than the bound $\Omega((1/\epsilon)^d)$ that would be
obtained by ignoring the geometric structure. This is the mathematical
formalization of the curse of dimensionality alleviated by the manifold
hypothesis.
\end{remark}

%=======================================================================
\section{Isometries Between Models of Hyperbolic Space}
\label{sec:app:hyperbolic-isometries}
%=======================================================================

There are three models of hyperbolic space (hyperboloid, Poincar\'e disc, upper
half-plane), and they are isometric as Riemannian manifolds. We prove the
isometries between them explicitly via stereographic projection and the Cayley
transform.

\subsection{Preliminaries}
\label{subsec:app:hyperbolic-prelim}

\begin{definition}[Riemannian isometry]
\label{def:app:riemannian-isometry}
Let $(\mathcal{M}, g)$ and $(\mathcal{N}, h)$ be Riemannian manifolds. A
diffeomorphism $F \colon \mathcal{M} \to \mathcal{N}$ is a \emph{Riemannian
isometry} if it preserves the metric, that is, if $F^*h = g$, where the
\emph{pullback} of the metric is defined by
\begin{equation}
\label{eq:app:metric-pullback}
(F^*h)_p(v, w) = h_{F(p)}(dF_p(v), dF_p(w))
\end{equation}
for all $p \in \mathcal{M}$ and $v, w \in T_p\mathcal{M}$.
\end{definition}

\begin{remark}
\label{rem:app:conformal-isometry}
When both metrics are \emph{conformal}, that is, of the form
$g_{ij} = \lambda(x)^2 \delta_{ij}$ and $h_{ij} = \mu(y)^2 \delta_{ij}$, checking
the condition $F^*h = g$ simplifies: it suffices to verify that
\[
\mu(F(x))^2 \cdot \norm{dF_x(v)}^2 = \lambda(x)^2 \cdot \norm{v}^2
\]
for every tangent vector $v$, where $\norm{\cdot}$ is the Euclidean norm. This
reduces the verification to a scalar identity between the conformal factors. See
\citet[Chapter~3]{lee2018introduction}.
\end{remark}

\subsection{From the hyperboloid to the Poincar\'e disc}
\label{subsec:app:hyperboloid-poincare}

\begin{theorem}[Hyperboloid--Poincar\'e disc isometry]
\label{thm:app:hyperboloid-poincare}
Stereographic projection from the south pole of the hyperboloid,
\begin{equation}
\label{eq:app:stereographic-hyperboloid}
\pi \colon \mathbb{H}^n \to \mathbb{D}^n , \qquad
\pi(x_0, x_1, \ldots, x_n) = \frac{(x_1, \ldots, x_n)}{1 + x_0} ,
\end{equation}
is a Riemannian isometry
$(\mathbb{H}^n, g_{\mathcal{L}}) \to (\mathbb{D}^n, g_{\mathbb{D}})$, where
$g_{\mathcal{L}}$ is the metric induced by the Lorentz product on the hyperboloid
and $g_{\mathbb{D}} = \frac{4}{(1-\norm{\mathbf{y}}^2)^2}\delta_{ij}$ is the
Poincar\'e metric. Its inverse is
\begin{equation}
\label{eq:app:inverse-stereographic}
\pi^{-1}(y_1, \ldots, y_n) = \frac{1}{1 - \norm{\mathbf{y}}^2}
\bigl(1 + \norm{\mathbf{y}}^2,\; 2y_1,\; \ldots,\; 2y_n\bigr) .
\end{equation}
\end{theorem}

\begin{proof}
We proceed in several steps.

\textbf{Step 1: Geometric interpretation.}
The map $\pi$ is the stereographic projection from the point
$(-1, 0, \ldots, 0)$ (the south pole of the hyperboloid in Minkowski space) onto
the hyperplane $\{x_0 = 0\} \cong \R^n$. The line joining $(-1, \mathbf{0})$ with
$(x_0, \mathbf{x}) \in \mathbb{H}^n$ is
$t \mapsto (-1, \mathbf{0}) + t\bigl((x_0+1), \mathbf{x}\bigr)$, and it meets the
hyperplane $x_0 = 0$ when $t = \frac{1}{1+x_0}$, giving the point
$\frac{\mathbf{x}}{1+x_0}$.

\textbf{Step 2: Bijectivity and smoothness.}
Both $\pi$ and $\pi^{-1}$ are rational functions of their arguments, hence smooth
on their domains. We verify that $\pi^{-1}$ is well defined: if
$\norm{\mathbf{y}} < 1$, we must check $\pi^{-1}(\mathbf{y}) \in \mathbb{H}^n$.
Let $s = \norm{\mathbf{y}}^2 < 1$. Then
\[
\inner{\pi^{-1}(\mathbf{y})}{\pi^{-1}(\mathbf{y})}_{\mathcal{L}}
= \frac{1}{(1-s)^2}\bigl(-(1+s)^2 + 4s\bigr)
= \frac{-(1-s)^2}{(1-s)^2} = -1 ,
\]
and the component $x_0 = \frac{1+s}{1-s} > 0$, so
$\pi^{-1}(\mathbf{y}) \in \mathbb{H}^n$.

For the composition $\pi \circ \pi^{-1}$: given $\mathbf{y} \in \mathbb{D}^n$,
$\pi^{-1}(\mathbf{y})$ has component
$x_0 = \frac{1+\norm{\mathbf{y}}^2}{1-\norm{\mathbf{y}}^2}$ and spatial components
$\frac{2y_i}{1-\norm{\mathbf{y}}^2}$. Then
\[
\pi(\pi^{-1}(\mathbf{y}))
= \frac{(2y_1, \ldots, 2y_n)/(1-\norm{\mathbf{y}}^2)}
{1 + (1+\norm{\mathbf{y}}^2)/(1-\norm{\mathbf{y}}^2)}
= \frac{(2y_1, \ldots, 2y_n)/(1-\norm{\mathbf{y}}^2)}
{2/(1-\norm{\mathbf{y}}^2)} = \mathbf{y} .
\]
The verification of $\pi^{-1} \circ \pi = \mathrm{id}$ is analogous.

\textbf{Step 3: Relation between norm and coordinates.}
Let $(x_0, \mathbf{x}) \in \mathbb{H}^n$ and
$\mathbf{y} = \pi(x_0, \mathbf{x}) = \frac{\mathbf{x}}{1+x_0}$. Using the
hyperboloid relation $x_0^2 - \norm{\mathbf{x}}^2 = 1$,
\begin{equation}
\label{eq:app:norm-y-hyperboloid}
\norm{\mathbf{y}}^2 = \frac{\norm{\mathbf{x}}^2}{(1+x_0)^2}
= \frac{x_0^2 - 1}{(1+x_0)^2} = \frac{x_0 - 1}{x_0 + 1} .
\end{equation}
This yields the key identity
\begin{equation}
\label{eq:app:conformal-factor-key}
1 - \norm{\mathbf{y}}^2 = 1 - \frac{x_0 - 1}{x_0 + 1} = \frac{2}{1 + x_0} .
\end{equation}

\textbf{Step 4: Computation of the differential.}
To compute $d\pi$, use intrinsic coordinates $(x_1, \ldots, x_n)$ on
$\mathbb{H}^n$. The component $x_0$ satisfies $x_0 = \sqrt{1 + \norm{\mathbf{x}}^2}$,
whence
\[
dx_0 = \frac{\sum_{i=1}^n x_i \, dx_i}{x_0} .
\]
Differentiating $y_j = \frac{x_j}{1+x_0}$,
\begin{align}
dy_j &= \frac{dx_j(1+x_0) - x_j \, dx_0}{(1+x_0)^2} \nonumber \\
&= \frac{dx_j}{1+x_0} - \frac{x_j}{(1+x_0)^2} \cdot
\frac{\sum_i x_i \, dx_i}{x_0} . \label{eq:app:differential-pi}
\end{align}

\textbf{Step 5: Verification of the isometry.}
We must check that $\pi^* g_{\mathbb{D}} = g_{\mathcal{L}}$. The Lorentz-induced
metric on $\mathbb{H}^n$ with intrinsic coordinates $(x_1, \ldots, x_n)$ is
\[
g_{\mathcal{L}} = -dx_0^2 + \sum_{i=1}^n dx_i^2
= -\frac{(\sum_i x_i \, dx_i)^2}{x_0^2} + \sum_{i=1}^n dx_i^2 ,
\]
that is, in components $(g_{\mathcal{L}})_{ij} = \delta_{ij} - \frac{x_i x_j}{x_0^2}$.
The pullback of the Poincar\'e metric is
\[
(\pi^* g_{\mathbb{D}})_{ij} = \frac{4}{(1-\norm{\mathbf{y}}^2)^2}
\sum_{k=1}^n \frac{\partial y_k}{\partial x_i} \frac{\partial y_k}{\partial x_j} .
\]
Using~\eqref{eq:app:conformal-factor-key}, the conformal factor is
$\frac{4}{(1-\norm{\mathbf{y}}^2)^2} = (1+x_0)^2$. From~\eqref{eq:app:differential-pi},
\[
\frac{\partial y_k}{\partial x_i} = \frac{\delta_{ki}}{1+x_0}
- \frac{x_k x_i}{x_0(1+x_0)^2} ,
\]
so that
\begin{align*}
\sum_{k} \frac{\partial y_k}{\partial x_i}\frac{\partial y_k}{\partial x_j}
&= \frac{\delta_{ij}}{(1+x_0)^2}
- \frac{2 x_i x_j}{x_0(1+x_0)^3}
+ \frac{x_i x_j (x_0^2 - 1)}{x_0^2(1+x_0)^4} \\
&= \frac{\delta_{ij}}{(1+x_0)^2}
- \frac{2 x_i x_j}{x_0(1+x_0)^3}
+ \frac{x_i x_j (x_0 - 1)}{x_0^2(1+x_0)^3} \\
&= \frac{\delta_{ij}}{(1+x_0)^2} - \frac{x_i x_j}{x_0^2(1+x_0)^2} ,
\end{align*}
where the last step uses
\[
-\frac{2}{x_0(1+x_0)} + \frac{x_0-1}{x_0^2(1+x_0)}
= \frac{-2x_0 + x_0 - 1}{x_0^2(1+x_0)}
= \frac{-(x_0+1)}{x_0^2(1+x_0)} = \frac{-1}{x_0^2} .
\]
Multiplying by $(1+x_0)^2$,
\[
(\pi^* g_{\mathbb{D}})_{ij} = (1+x_0)^2 \left(\frac{\delta_{ij}}{(1+x_0)^2}
- \frac{x_i x_j}{x_0^2(1+x_0)^2}\right)
= \delta_{ij} - \frac{x_i x_j}{x_0^2} = (g_{\mathcal{L}})_{ij} .
\]
This shows $\pi^* g_{\mathbb{D}} = g_{\mathcal{L}}$, that is, $\pi$ is an
isometry.
\end{proof}

\subsection{From the Poincar\'e disc to the upper half-plane}
\label{subsec:app:poincare-halfplane}

\begin{theorem}[Poincar\'e disc--upper half-plane isometry]
\label{thm:app:poincare-halfplane}
In dimension $n = 2$, the Cayley transform
\begin{equation}
\label{eq:app:cayley-transform}
C \colon \mathbb{D}^2 \to \mathbb{H}^2 , \qquad C(w) = i\,\frac{1+w}{1-w} ,
\end{equation}
is a Riemannian isometry from the Poincar\'e disc
$(\mathbb{D}^2, g_{\mathbb{D}})$ to the upper half-plane
$(\mathbb{H}^2, g_{\mathbb{H}})$, where $g_{\mathbb{H}} = \frac{dx^2 + dy^2}{y^2}$.
Its inverse is
\begin{equation}
\label{eq:app:cayley-inverse}
C^{-1} \colon \mathbb{H}^2 \to \mathbb{D}^2 , \qquad C^{-1}(z) = \frac{z - i}{z + i} .
\end{equation}
\end{theorem}

\begin{proof}
We proceed in four steps.

\textbf{Step 1: Well-definedness ($C(\mathbb{D}^2) \subseteq \mathbb{H}^2$).}
Let $w \in \mathbb{D}^2$, that is, $\abs{w} < 1$. We must verify
$\Img(C(w)) > 0$. A direct computation gives
\[
\frac{1+w}{1-w} = \frac{(1+w)(1-\bar{w})}{\abs{1-w}^2}
= \frac{1 - \bar{w} + w - \abs{w}^2}{\abs{1-w}^2} ,
\]
hence
\[
\Img(C(w)) = \operatorname{Re}\!\left(\frac{1+w}{1-w}\right)
= \frac{1 - \abs{w}^2}{\abs{1-w}^2} .
\]
Since $\abs{w} < 1$, we have $1 - \abs{w}^2 > 0$ and $\abs{1-w}^2 > 0$, so
$\Img(C(w)) > 0$.

\textbf{Step 2: Bijectivity.}
For $C^{-1}(z) = \frac{z-i}{z+i}$ with $\Img(z) > 0$, writing $z = a + bi$ with
$b > 0$,
\[
\abs{C^{-1}(z)}^2 = \frac{\abs{z-i}^2}{\abs{z+i}^2}
= \frac{a^2 + (b-1)^2}{a^2 + (b+1)^2}
= \frac{a^2 + b^2 - 2b + 1}{a^2 + b^2 + 2b + 1} < 1 ,
\]
since $-2b < 2b$ when $b > 0$. Hence $C^{-1}(z) \in \mathbb{D}^2$. The
compositions $C \circ C^{-1} = \mathrm{id}$ and $C^{-1} \circ C = \mathrm{id}$
are verified by direct computation.

\textbf{Step 3: Differential.}
The complex derivative of $C$ is
\begin{equation}
\label{eq:app:cayley-derivative}
C'(w) = i \cdot \frac{(1-w) - (1+w)(-1)}{(1-w)^2} = \frac{2i}{(1-w)^2} .
\end{equation}

\textbf{Step 4: Verification of the isometry.}
For a holomorphic map $C$ between surfaces with conformal metrics, the pullback is
\[
(C^* g_{\mathbb{H}})_w = \frac{\abs{C'(w)}^2}{(\Img C(w))^2} \cdot (du^2 + dv^2) ,
\]
where $w = u + iv$. Substituting the results of the previous steps,
\[
\frac{\abs{C'(w)}^2}{(\Img C(w))^2}
= \frac{4/\abs{1-w}^4}{(1-\abs{w}^2)^2/\abs{1-w}^4}
= \frac{4}{(1-\abs{w}^2)^2} ,
\]
which is exactly the conformal factor of the Poincar\'e metric
$g_{\mathbb{D}} = \frac{4}{(1-\abs{w}^2)^2}(du^2 + dv^2)$. Therefore
$C^* g_{\mathbb{H}} = g_{\mathbb{D}}$.
\end{proof}

\subsection{Equivalence of the three models}
\label{subsec:app:three-models}

\begin{corollary}[Equivalence of the three hyperbolic models]
\label{cor:app:hyperbolic-equivalence}
The three models of hyperbolic space (hyperboloid, Poincar\'e disc, and upper
half-plane) are Riemannian isometric. In particular, the composition
$C \circ \pi$ provides an explicit isometry from the hyperboloid model to the
upper half-plane.
\end{corollary}

\begin{proof}
A composition of the isometries of
\Cref{thm:app:hyperboloid-poincare,thm:app:poincare-halfplane}: since the
composition of isometries is an isometry,
$C \circ \pi \colon (\mathbb{H}^n, g_{\mathcal{L}}) \to
(\mathbb{D}^n, g_{\mathbb{D}}) \to (\mathbb{H}^2, g_{\mathbb{H}})$ is an isometry
(in dimension $n = 2$ for the half-plane).
\end{proof}

\begin{remark}[Implications for geometric deep learning]
\label{rem:app:models-gdl}
The isometric equivalence of the three models justifies treating the choice of
model in GDL applications as a matter of computational convenience:
\begin{itemize}
    \item The \textbf{Poincar\'e disc} $\mathbb{D}^n$ is preferred for
          gradient-based optimization, since the exponential and logarithmic maps
          admit closed-form expressions~\citep{nickel2017poincare}.
    \item The \textbf{hyperboloid model} offers greater numerical stability,
          since the conformal factor $\frac{4}{(1-\norm{\mathbf{x}}^2)^2}$ of the
          Poincar\'e disc diverges as $\norm{\mathbf{x}} \to 1$ (near the boundary
          of the disc), whereas the induced Lorentz metric has no such
          singularity.
    \item The transformation $\pi$ of \Cref{thm:app:hyperboloid-poincare} is
          precisely the map used in practice to convert between representations.
\end{itemize}
\end{remark}

%=======================================================================
\section{Sectional Curvature of the Model Spaces}
\label{sec:app:model-curvature}
%=======================================================================

The sphere $S^n$ has constant sectional curvature $+1$ and hyperbolic space
$\mathbb{H}^n$ has constant sectional curvature $-1$. We prove both results by
applying the definition of sectional curvature to the explicit metrics of these
spaces.

\subsection{Preliminaries}
\label{subsec:app:model-curvature-prelim}

\begin{lemma}[Gauss equation for hypersurfaces]
\label{lem:app:gauss-equation}
Let $M^n \subset \R^{n+1}$ be an oriented hypersurface with unit normal $\nu$,
and let $S_p \colon T_pM \to T_pM$ be the \emph{shape operator} (Weingarten map)
$S_p(v) = -d\nu_p(v)$. The second fundamental form is
$\II(u, v) = \inner{S_p(u)}{v}$. If $u, v \in T_pM$ are orthonormal, then the
sectional curvature of $M$ in the plane spanned by $\{u, v\}$ satisfies
\begin{equation}
\label{eq:app:gauss-equation}
K_M(u,v) = K_{\R^{n+1}}(u,v) + \II(u,u)\,\II(v,v) - \II(u,v)^2 .
\end{equation}
Since $\R^{n+1}$ is flat ($K_{\R^{n+1}} \equiv 0$), this reduces to
\begin{equation}
\label{eq:app:gauss-simplified}
K_M(u,v) = \II(u,u)\,\II(v,v) - \II(u,v)^2 .
\end{equation}
\end{lemma}

The proof can be found in \citet[Theorem~8.4]{lee2018introduction} or in
\citet[Chapter~3, Proposition~4]{do1976differential}.

\begin{lemma}[Curvature of conformal metrics in dimension 2]
\label{lem:app:conformal-curvature-2d}
Let $g = \lambda(x,y)^2(dx^2 + dy^2)$ be a conformal metric on a domain of
$\R^2$, where $\lambda > 0$ is smooth. Then the Gaussian curvature of
$(\R^2, g)$ is
\begin{equation}
\label{eq:app:conformal-curvature}
K = -\frac{\Delta(\ln \lambda)}{\lambda^2} ,
\end{equation}
where $\Delta = \partial_{xx} + \partial_{yy}$ is the Euclidean Laplacian.
\end{lemma}

\begin{proof}
For the metric $g_{ij} = \lambda^2 \delta_{ij}$, the Christoffel symbols are
\[
\Gamma^k_{ij} = (\partial_i \ln\lambda)\,\delta_{jk}
+ (\partial_j \ln\lambda)\,\delta_{ik} - (\partial_k \ln\lambda)\,\delta_{ij} .
\]
Substituting into the expression for the Riemann tensor $R^l{}_{ijk}$ and
contracting to obtain the Gaussian curvature $K = R^1{}_{212}/\lambda^2$, a direct
computation gives~\eqref{eq:app:conformal-curvature}. The details can be found in
\citet[Proposition~8.32]{lee2018introduction}.
\end{proof}

\subsection{Curvature of the sphere \texorpdfstring{$S^n$}{Sn}}
\label{subsec:app:sphere-curvature}

\begin{theorem}[Sectional curvature of the sphere]
\label{thm:app:sphere-curvature}
The unit sphere $S^n \subset \R^{n+1}$ with the Riemannian metric induced by the
ambient Euclidean metric has constant sectional curvature $K = +1$.
\end{theorem}

\begin{proof}
Consider $S^n = \{p \in \R^{n+1} : \norm{p} = 1\}$ with the metric induced by
$\R^{n+1}$.

\textbf{Step 1: Unit normal.}
The outward unit normal field is $\nu(p) = p$ for all $p \in S^n$.

\textbf{Step 2: Shape operator.}
For $v \in T_pS^n$, consider the curve
$\gamma(t) = \frac{p + tv}{\norm{p + tv}}$ on $S^n$ with $\gamma(0) = p$ and
$\gamma'(0) = v$ (since $\inner{p}{v} = 0$). Then
$\nu(\gamma(t)) = \gamma(t) = \frac{p + tv}{\norm{p + tv}}$, and differentiating
at $t = 0$,
\[
d\nu_p(v) = \frac{d}{dt}\bigg|_{t=0} \frac{p + tv}{\norm{p + tv}} = v ,
\]
since $\norm{p + tv}^2 = 1 + t^2\norm{v}^2$ implies
$\frac{d}{dt}\norm{p+tv}\big|_{t=0} = 0$. Hence the shape operator is
$S_p = -d\nu_p = -\operatorname{Id}_{T_pS^n}$.

\textbf{Step 3: Second fundamental form.}
For orthonormal $u, v \in T_pS^n$,
\[
\II(u,u) = \inner{S_p(u)}{u} = \inner{-u}{u} = -1 , \quad
\II(v,v) = -1 , \quad \II(u,v) = \inner{-u}{v} = 0 .
\]

\textbf{Step 4: Gauss equation.}
Applying \Cref{lem:app:gauss-equation},
\[
K(u,v) = (-1)(-1) - 0^2 = 1 .
\]
Since the choice of $p \in S^n$ and of the $2$-plane $\{u,v\}$ is arbitrary,
$K \equiv +1$.
\end{proof}

\subsection{Curvature of hyperbolic space \texorpdfstring{$\mathbb{H}^n$}{Hn}}
\label{subsec:app:hyperbolic-curvature}

\begin{theorem}[Sectional curvature of hyperbolic space]
\label{thm:app:hyperbolic-curvature}
The Poincar\'e disc $(\mathbb{D}^n, g_{\mathbb{D}})$ with metric
$g_{\mathbb{D}} = \frac{4}{(1-\norm{\mathbf{x}}^2)^2}\delta_{ij}$ has constant
sectional curvature $K = -1$.
\end{theorem}

\begin{proof}
We proceed in two parts: first an explicit computation of the curvature in
dimension $2$, then an extension to arbitrary dimension.

\textbf{Part 1: Dimension 2 (explicit computation).}
On $\mathbb{D}^2$ the metric is conformal with factor
$\lambda(x,y) = \frac{2}{1 - r^2}$, where $r^2 = x^2 + y^2$. We compute
$\ln\lambda = \ln 2 - \ln(1 - r^2)$ and
\[
\partial_x(\ln\lambda) = \frac{2x}{1 - r^2} , \qquad
\partial_{xx}(\ln\lambda) = \frac{2(1-r^2) + 4x^2}{(1-r^2)^2} .
\]
By symmetry, $\partial_{yy}(\ln\lambda) = \frac{2(1-r^2) + 4y^2}{(1-r^2)^2}$.
Summing,
\[
\Delta(\ln\lambda) = \frac{4(1-r^2) + 4r^2}{(1-r^2)^2} = \frac{4}{(1-r^2)^2} .
\]
Applying \Cref{lem:app:conformal-curvature-2d},
\[
K = -\frac{\Delta(\ln\lambda)}{\lambda^2}
= -\frac{4/(1-r^2)^2}{4/(1-r^2)^2} = -1 .
\]

\textbf{Part 2: Extension to dimension $n$ (symmetry argument).}
The metric $g_{\mathbb{D}}$ is invariant under the action of $\Orth(n)$ by
rotations of $\mathbf{x}$ (since $g_{\mathbb{D}}$ depends only on
$\norm{\mathbf{x}}$). At any point $p \in \mathbb{D}^n$, the stabilizer of $p$
under this action acts transitively on the $2$-planes of $T_p\mathbb{D}^n$ that
contain the radial direction, and by composition with rotations, on all
$2$-planes. Hence all sectional curvatures at $p$ are equal.

The subspace $\Pi = \{x \in \mathbb{D}^n : x_3 = \cdots = x_n = 0\}$ is a totally
geodesic submanifold isometric to $(\mathbb{D}^2, g_{\mathbb{D}^2})$: the
reflection $\sigma \colon x_i \mapsto -x_i$ for $i \geq 3$ is an isometry of
$(\mathbb{D}^n, g_{\mathbb{D}})$ whose fixed-point set is $\Pi$, and the
fixed-point set of an isometry is totally geodesic
\citep[Proposition~5.21]{lee2018introduction}. By Part~1, the sectional curvature
on $\Pi$ is $-1$, so $K = -1$ for all $2$-planes at each point of $\Pi$.

Finally, the Poincar\'e metric is homogeneous: M\"obius transformations provide
isometries carrying any point of $\mathbb{D}^n$ to any other
\citep{nickel2017poincare}. Therefore $K = -1$ throughout $\mathbb{D}^n$.
\end{proof}

\begin{corollary}
\label{cor:app:hyperbolic-models-curvature}
By \Cref{thm:app:hyperboloid-poincare,thm:app:poincare-halfplane}, the
hyperboloid model $(\mathbb{H}^n, g_{\mathcal{L}})$ and the upper half-plane
$(\mathbb{H}^2, g_{\mathbb{H}})$ also have constant sectional curvature
$K = -1$.
\end{corollary}

\begin{proof}
Isometries preserve sectional curvature: if $F \colon (M, g) \to (N, h)$ is a
Riemannian isometry, then $K_M(\sigma) = K_N(dF(\sigma))$ for every $2$-plane
$\sigma \subset T_pM$. The result follows from
\Cref{thm:app:hyperbolic-curvature} and the explicit isometries of
\Cref{sec:app:hyperbolic-isometries}.
\end{proof}

%=======================================================================
\section{Lagrange's Theorem}
\label{sec:app:lagrange}
%=======================================================================

Lagrange's theorem is a fundamental result of finite group theory relating the
order of a group to the orders of its subgroups. Although its proof is
elementary, we include the details for completeness.

\begin{theorem}[Lagrange]
\label{thm:app:lagrange}
Let $\group$ be a finite group and $H \leq \group$ a subgroup. Then $\abs{H}$
divides $\abs{\group}$, and
\begin{equation}
\label{eq:app:lagrange}
\abs{\group} = \abs{H} \cdot [\group:H] ,
\end{equation}
where $[\group:H]$ denotes the index of $H$ in $\group$ (the number of left
cosets of $H$ in $\group$).
\end{theorem}

\begin{proof}
Define an equivalence relation on $\group$ by $a \sim b$ if and only if
$a^{-1}b \in H$. The equivalence classes are the left cosets
$g H = \{gh : h \in H\}$ for $g \in \group$.

\textbf{Step 1.} Each coset has exactly $\abs{H}$ elements. Indeed, the map
$\varphi_g \colon H \to gH$ given by $\varphi_g(h) = gh$ is bijective: surjective
by construction, and injective because $gh_1 = gh_2$ implies $h_1 = h_2$ by the
cancellation law.

\textbf{Step 2.} The cosets partition $\group$. This is a direct consequence of
$\sim$ being an equivalence relation.

\textbf{Step 3.} If $[\group:H]$ denotes the number of distinct cosets, then
\[
\abs{\group} = [\group:H] \cdot \abs{H} ,
\]
since $\group$ decomposes into exactly $[\group:H]$ disjoint sets, each of
cardinality $\abs{H}$. In particular, $\abs{H}$ divides $\abs{\group}$.
\end{proof}

%=======================================================================
\section{Rad\'o's Theorem (Triangulability of Surfaces)}
\label{sec:app:rado}
%=======================================================================

Rad\'o's theorem states that every compact surface admits a triangulation, a
result fundamental both for the proof of the Gauss--Bonnet theorem
(\Cref{thm:app:gauss-bonnet-general}) and for applications of geometric deep
learning on triangulated meshes.

\subsection{Preliminaries}
\label{subsec:app:rado-prelim}

Recall that a \emph{topological $2$-manifold} (or \emph{surface}) is a
topological space $M$ that is Hausdorff, second countable, and such that every
point $p \in M$ has a neighbourhood homeomorphic to an open subset of $\R^2$. The
following classical result of planar topology is a key tool in the proof.

\begin{theorem}[Jordan--Schoenflies theorem]
\label{thm:app:jordan-schoenflies}
Every homeomorphism $h \colon S^1 \to \R^2$ whose image is a Jordan curve
$\mathcal{C}$ extends to a homeomorphism
$H \colon \overline{D}^2 \to \overline{U}$, where $D^2$ is the open unit disc and
$U$ is the bounded component of $\R^2 \setminus \mathcal{C}$.
\end{theorem}

For a proof, see \citet[Chapter~4]{moise1977geometric}.

\subsection{Formal statement}
\label{subsec:app:rado-statement}

\begin{theorem}[Rad\'o's theorem, 1925]
\label{thm:app:rado}
Every compact topological $2$-manifold admits a triangulation. That is, if $M$
is a compact surface, there exist a simplicial complex $\mathcal{K}$ and a
homeomorphism $h \colon \abs{\mathcal{K}} \to M$.
\end{theorem}

\begin{remark}[Triangulability in higher dimensions]
\label{rem:app:triangulability-dimensions}
Triangulability of manifolds depends crucially on the dimension:
\begin{itemize}
    \item \textbf{Dimension 2}: Every compact surface is triangulable
          (\Cref{thm:app:rado}).
    \item \textbf{Dimension 3}: \citet{moise1977geometric} proved in 1952 that
          every compact $3$-manifold admits a triangulation, a considerably
          harder result.
    \item \textbf{Dimension $\geq 4$}: There exist non-triangulable compact
          topological manifolds. Casson's result (based on Freedman's work on
          $4$-manifolds) shows that the manifold $E_8$ admits no triangulation at
          all.
\end{itemize}
\end{remark}

\subsection{Proof sketch}
\label{subsec:app:rado-proof}

We present the main ideas of the proof, following the exposition of
\citet[Chapter~8]{moise1977geometric}. See also \citet[\S77]{munkres2000topology}
for an alternative treatment.

\begin{proof}[Proof sketch of \Cref{thm:app:rado}]
The proof proceeds in four steps.

\textbf{Step 1: Local triangulability.}
Since $M$ is a $2$-manifold, each point $p \in M$ has a neighbourhood $U_p$
homeomorphic to an open subset of $\R^2$ via a chart $(U_p, \phi_p)$. In $\R^2$,
every bounded region admits a triangulation (for example, by barycentric
subdivision of a grid). Hence each $U_p$ can be triangulated locally by
transporting this triangulation via $\phi_p^{-1}$.

\textbf{Step 2: Finite cover.}
By compactness of $M$, there is a finite subcover $\{U_1, \ldots, U_N\}$ with
local triangulations $\mathcal{T}_1, \ldots, \mathcal{T}_N$.

\textbf{Step 3: Compatibility via Jordan--Schoenflies.}
The essential step is to make the local triangulations compatible on the overlap
regions $U_i \cap U_j$. The \Cref{thm:app:jordan-schoenflies} allows one to
extend homeomorphisms between simple closed curves in $\R^2$ to homeomorphisms of
their interiors, which guarantees that the local triangulations can be adjusted to
be compatible on the intersections.

\textbf{Step 4: Refinement and gluing.}
By successive subdivisions (refinements) of the local triangulations, one obtains
triangulations fine enough to be compatible on all intersections. Gluing produces
a global simplicial complex $\mathcal{K}$ whose geometric realization
$\abs{\mathcal{K}}$ is homeomorphic to $M$.
\end{proof}

%=======================================================================
\section{The Spectral Theorem for Real Symmetric Matrices}
\label{sec:app:spectral}
%=======================================================================

The spectral theorem is a fundamental result of linear algebra guaranteeing the
orthogonal diagonalization of every real symmetric matrix. In the context of
geometric deep learning, this theorem underlies the spectral decomposition of the
graph Laplacian, the graph Fourier transform, and the definition of
positive-semidefinite matrices. For a detailed exposition, see
\citet[Chapter~7]{Axler2015}.

\begin{theorem}[Spectral theorem for real symmetric matrices]
\label{thm:app:spectral}
Let $A \in \R^{n \times n}$ be a symmetric matrix ($A = A^T$). Then:
\begin{enumerate}
    \item All eigenvalues of $A$ are real.
    \item Eigenvectors corresponding to distinct eigenvalues are orthogonal.
    \item There exists an orthogonal matrix $U \in \Orth(n)$ such that
          \[
          A = U \Lambda U^T , \quad
          \Lambda = \diag(\lambda_1, \ldots, \lambda_n) ,
          \]
          where $\lambda_1, \ldots, \lambda_n$ are the eigenvalues of $A$.
\end{enumerate}
\end{theorem}

\begin{proof}
We proceed by induction on $n$.

\textbf{Base case $n = 1$.} Every matrix $A = (a) \in \R^{1 \times 1}$ has a
single eigenvalue $\lambda_1 = a \in \R$, and $U = (1)$ satisfies
$A = U \Lambda U^T$ trivially.

\textbf{Inductive step.} Assume the result holds for symmetric matrices of size
$(n-1) \times (n-1)$ and let $A \in \R^{n \times n}$ be symmetric.

\textbf{(1) Real eigenvalues.} Let $\lambda$ be an eigenvalue of $A$ (possibly
complex) with eigenvector $\mathbf{v} \in \C^n \setminus \{\mathbf{0}\}$. Using
the standard Hermitian product
$\inner{\mathbf{u}}{\mathbf{w}}_{\C} = \overline{\mathbf{u}}^T \mathbf{w}$,
\[
\lambda \inner{\mathbf{v}}{\mathbf{v}}_{\C}
= \inner{\mathbf{v}}{A\mathbf{v}}_{\C}
= \inner{A^T \mathbf{v}}{\mathbf{v}}_{\C}
= \inner{A\mathbf{v}}{\mathbf{v}}_{\C}
= \overline{\inner{\mathbf{v}}{A\mathbf{v}}_{\C}}
= \bar{\lambda} \inner{\mathbf{v}}{\mathbf{v}}_{\C} ,
\]
where we used that $A = A^T$ implies $A = \bar{A}^T = A^*$ (since $A$ is real).
As $\inner{\mathbf{v}}{\mathbf{v}}_{\C} = \norm{\mathbf{v}}^2 > 0$, we get
$\lambda = \bar{\lambda}$, that is, $\lambda \in \R$.

\textbf{(2) Orthogonality of eigenvectors.} Let $A\mathbf{v} = \lambda \mathbf{v}$
and $A\mathbf{w} = \mu \mathbf{w}$ with $\lambda \neq \mu$ (both real by (1), and
$\mathbf{v}, \mathbf{w} \in \R^n$). Then
\[
\lambda \inner{\mathbf{v}}{\mathbf{w}} = \inner{A\mathbf{v}}{\mathbf{w}}
= \inner{\mathbf{v}}{A\mathbf{w}} = \mu \inner{\mathbf{v}}{\mathbf{w}} ,
\]
where $\inner{\cdot}{\cdot}$ is the Euclidean inner product on $\R^n$ and the
second equality uses the symmetry of $A$. Hence
$(\lambda - \mu)\inner{\mathbf{v}}{\mathbf{w}} = 0$, and since
$\lambda \neq \mu$, we conclude $\inner{\mathbf{v}}{\mathbf{w}} = 0$.

\textbf{(3) Orthogonal diagonalization.} By (1), $A$ has a real eigenvalue
$\lambda_1$ with unit eigenvector $\mathbf{v}_1 \in \R^n$ ($\norm{\mathbf{v}_1} = 1$).
Consider the orthogonal complement
$W = \mathbf{v}_1^\perp = \{\mathbf{w} \in \R^n : \inner{\mathbf{w}}{\mathbf{v}_1}
= 0\}$, a subspace of dimension $n - 1$.

\textit{Claim}: $W$ is $A$-invariant, that is, $A\mathbf{w} \in W$ for all
$\mathbf{w} \in W$. Indeed, if $\mathbf{w} \in W$,
\[
\inner{A\mathbf{w}}{\mathbf{v}_1} = \inner{\mathbf{w}}{A\mathbf{v}_1}
= \lambda_1 \inner{\mathbf{w}}{\mathbf{v}_1} = 0 ,
\]
so $A\mathbf{w} \in \mathbf{v}_1^\perp = W$.

Let $\{\mathbf{e}_2, \ldots, \mathbf{e}_n\}$ be an orthonormal basis of $W$ and
$P = [\mathbf{e}_2 \mid \cdots \mid \mathbf{e}_n] \in \R^{n \times (n-1)}$. The
restriction $A|_W$ is represented by the symmetric matrix
$B = P^T A P \in \R^{(n-1) \times (n-1)}$. By the inductive hypothesis, there is
an orthogonal matrix $\tilde{U} \in \Orth(n-1)$ with
$B = \tilde{U} \tilde{\Lambda} \tilde{U}^T$, where
$\tilde{\Lambda} = \diag(\lambda_2, \ldots, \lambda_n)$. Define
\[
U = \begin{pmatrix} \mathbf{v}_1 & P\tilde{U} \end{pmatrix} \in \R^{n \times n} .
\]
Then $U$ is orthogonal (its columns form an orthonormal basis of $\R^n$) and
\[
U^T A U = \diag(\lambda_1, \lambda_2, \ldots, \lambda_n) ,
\]
which is equivalent to $A = U \Lambda U^T$, completing the induction.
\end{proof}

\begin{corollary}[Characterization of positive-semidefinite matrices]
\label{cor:app:psd-characterization}
Let $A \in \R^{n \times n}$ be symmetric with eigenvalues
$\lambda_1, \ldots, \lambda_n$. Then $A$ is positive semidefinite
($A \succeq 0$) if and only if $\lambda_i \geq 0$ for all
$i \in \{1, \ldots, n\}$.
\end{corollary}

\begin{proof}
By \Cref{thm:app:spectral}, $A = U \Lambda U^T$ with $U$ orthogonal and
$\Lambda = \diag(\lambda_1, \ldots, \lambda_n)$. For any $\mathbf{x} \in \R^n$,
let $\mathbf{y} = U^T \mathbf{x}$. Then
\[
\mathbf{x}^T A \mathbf{x} = \mathbf{x}^T U \Lambda U^T \mathbf{x}
= \mathbf{y}^T \Lambda \mathbf{y} = \sum_{i=1}^n \lambda_i y_i^2 .
\]
Since $U$ is orthogonal, the map $\mathbf{x} \mapsto \mathbf{y} = U^T \mathbf{x}$
is a bijection of $\R^n$. If all $\lambda_i \geq 0$, the sum is clearly $\geq 0$
for every $\mathbf{x}$. Conversely, if some $\lambda_j < 0$, taking
$\mathbf{x} = U\mathbf{e}_j$ (where $\mathbf{e}_j$ is the $j$-th standard basis
vector) gives $\mathbf{x}^T A \mathbf{x} = \lambda_j < 0$.
\end{proof}

%=======================================================================
\section{The Inverse Function Theorem}
\label{sec:app:inverse-function}
%=======================================================================

The inverse function theorem is a fundamental tool of analysis that guarantees the
local existence of differentiable inverses. In the context of smooth manifolds it
is used repeatedly to construct local diffeomorphisms, in particular to show that
the Riemannian exponential map is a local diffeomorphism.

\subsection{Preliminary: the Banach fixed-point theorem}
\label{subsec:app:banach-fixed-point}

\begin{theorem}[Banach fixed-point theorem]
\label{thm:app:banach-fixed-point}
Let $(X, d)$ be a complete metric space and $T \colon X \to X$ a
\emph{contraction}, that is, there exists $\lambda \in [0, 1)$ such that
\[
d(T(x), T(y)) \leq \lambda \, d(x, y) \quad \text{for all } x, y \in X .
\]
Then $T$ has a unique fixed point $x^* \in X$ (i.e.\ $T(x^*) = x^*$). Moreover,
for any $x_0 \in X$, the sequence $x_{n+1} = T(x_n)$ converges to $x^*$, with the
estimate
\[
d(x_n, x^*) \leq \frac{\lambda^n}{1 - \lambda} \, d(x_0, x_1) .
\]
\end{theorem}

\begin{proof}
\textbf{Existence.} Let $x_0 \in X$ be arbitrary and set $x_{n+1} = T(x_n)$. For
$n \geq 1$,
\[
d(x_{n+1}, x_n) = d(T(x_n), T(x_{n-1})) \leq \lambda \, d(x_n, x_{n-1})
\leq \cdots \leq \lambda^n \, d(x_1, x_0) .
\]
For $m > n$, the triangle inequality gives
\[
d(x_m, x_n) \leq \sum_{k=n}^{m-1} d(x_{k+1}, x_k)
\leq \sum_{k=n}^{m-1} \lambda^k \, d(x_1, x_0)
\leq \frac{\lambda^n}{1 - \lambda} \, d(x_1, x_0) .
\]
Since $\lambda < 1$, the right-hand side tends to $0$ as $n \to \infty$, so
$(x_n)$ is a Cauchy sequence. By completeness of $(X, d)$, there is a limit
$x^* = \lim_{n \to \infty} x_n$.

By continuity of $T$ (every contraction is Lipschitz, hence continuous),
\[
T(x^*) = T\!\left(\lim_{n \to \infty} x_n\right)
= \lim_{n \to \infty} T(x_n) = \lim_{n \to \infty} x_{n+1} = x^* .
\]

\textbf{Uniqueness.} If $T(y^*) = y^*$ for another point $y^* \in X$, then
\[
d(x^*, y^*) = d(T(x^*), T(y^*)) \leq \lambda \, d(x^*, y^*) .
\]
Since $\lambda < 1$, this forces $d(x^*, y^*) = 0$, that is, $x^* = y^*$.

\textbf{Estimate.} Letting $m \to \infty$ in the inequality
$d(x_m, x_n) \leq \frac{\lambda^n}{1-\lambda} d(x_1, x_0)$ yields the result.
\end{proof}

\subsection{Statement and proof of the inverse function theorem}
\label{subsec:app:inverse-function-proof}

\begin{theorem}[Inverse function theorem]
\label{thm:app:inverse-function}
Let $U \subseteq \R^n$ be open, $F \colon U \to \R^n$ a map of class $C^k$
($k \geq 1$), and $p \in U$ a point such that the differential
$DF_p \colon \R^n \to \R^n$ is invertible. Then there exist open neighbourhoods
$V$ of $p$ and $W$ of $F(p)$ such that $F|_V \colon V \to W$ is a $C^k$
diffeomorphism.
\end{theorem}

\begin{proof}
The proof follows the classical strategy of \citet[Theorem~9.24]{Rudin1976},
presented in four steps.

\textbf{Step 1: Reduction to the case $DF_p = I$.}
Consider $G = (DF_p)^{-1} \circ F$. Then $G$ is of class $C^k$,
$G(p) = (DF_p)^{-1}(F(p))$, and $DG_p = (DF_p)^{-1} \circ DF_p = I$. If we prove
the theorem for $G$, composing with the invertible linear map $DF_p$ (a global
diffeomorphism) yields the result for $F$. Thus we may assume without loss of
generality that $DF_p = I$ and, after translating, that $p = 0$ and $F(0) = 0$.

\textbf{Step 2: Local surjectivity via the Banach fixed point.}
Define $\varphi(x) = x - F(x)$, so that $D\varphi_0 = I - DF_0 = 0$. By continuity
of $D\varphi$, there is $r > 0$ with $\norm{D\varphi_x} \leq \tfrac{1}{2}$ for all
$x \in \overline{B}(0, r)$. By the mean value theorem, for
$x_1, x_2 \in \overline{B}(0, r)$,
\[
\norm{\varphi(x_1) - \varphi(x_2)} \leq \tfrac{1}{2} \norm{x_1 - x_2} .
\]
For each $y \in B(0, r/2)$, define
$T_y \colon \overline{B}(0, r) \to \R^n$ by
$T_y(x) = y + \varphi(x) = y + x - F(x)$. Then
\[
\norm{T_y(x)} \leq \norm{y} + \norm{\varphi(x) - \varphi(0)}
\leq \tfrac{r}{2} + \tfrac{1}{2}\norm{x} \leq r ,
\]
so $T_y$ maps $\overline{B}(0, r)$ into itself. Moreover, $T_y$ is a contraction
with constant $\tfrac{1}{2}$, since
$T_y(x_1) - T_y(x_2) = \varphi(x_1) - \varphi(x_2)$. As $\overline{B}(0, r)$ is
complete, \Cref{thm:app:banach-fixed-point} guarantees a unique
$x \in \overline{B}(0, r)$ with $T_y(x) = x$, equivalently $F(x) = y$. This shows
that $F$ is surjective onto $W = B(0, r/2)$.

\textbf{Step 3: Injectivity.}
If $F(x_1) = F(x_2)$ with $x_1, x_2 \in \overline{B}(0, r)$, then
\[
\norm{x_1 - x_2} = \norm{\varphi(x_1) - \varphi(x_2) + F(x_1) - F(x_2)}
= \norm{\varphi(x_1) - \varphi(x_2)} \leq \tfrac{1}{2}\norm{x_1 - x_2} ,
\]
whence $x_1 = x_2$. Let $V = F^{-1}(W) \cap B(0, r)$; then
$F|_V \colon V \to W$ is bijective.

\textbf{Step 4: Differentiability of $F^{-1}$.}
Let $g = (F|_V)^{-1} \colon W \to V$. We must show $g$ is of class $C^k$. For
$y, y + k \in W$ with $x = g(y)$ and $x + h = g(y + k)$,
\[
k = F(x + h) - F(x) = DF_x(h) + \norm{h} \cdot \varepsilon(h) ,
\]
where $\varepsilon(h) \to 0$ as $h \to 0$. From the contraction estimate,
$\norm{h} \leq 2\norm{k}$, so $h \to 0$ as $k \to 0$. Since $DF_x$ is invertible
(its determinant is a continuous nonzero function),
\[
g(y + k) - g(y) = h = (DF_x)^{-1}(k) + \norm{k} \cdot \tilde{\varepsilon}(k) ,
\]
with $\tilde{\varepsilon}(k) \to 0$, so $Dg_y = (DF_{g(y)})^{-1}$. Since matrix
inversion and composition are $C^\infty$ operations, an inductive argument shows
that $g$ is of class $C^k$.
\end{proof}

%=======================================================================
\section{Preliminaries for the Universal Approximation Theorem}
\label{sec:app:uat-preliminaries}
%=======================================================================

The universal approximation theorem of \Cref{ch:foundations} guarantees that any
continuous function can be approximated arbitrarily well by an artificial neural
network with a single hidden layer. Its proof requires tools from measure theory
and functional analysis: the notions of measure and $\sigma$-algebra underpin the
representation (Riesz) and separation (Hahn--Banach) theorems that form the core
of the proof. We collect here the necessary preliminaries, leaving the statement
and proof of the theorem itself to \Cref{ch:foundations}. For an in-depth
treatment of these foundations, see \citet{Folland1999} or \citet{Rudin1976}.

\subsection{Measure-theoretic preliminaries}
\label{subsec:app:measure-theory}

\begin{definition}[$\sigma$-algebra]
\label{def:app:sigma-algebra}
Let $X$ be a nonempty set. A \emph{$\sigma$-algebra} (or $\sigma$-field) on $X$ is
a collection $\mathcal{F}$ of subsets of $X$ such that:
\begin{enumerate}
    \item $X \in \mathcal{F}$;
    \item if $A \in \mathcal{F}$, then $A^c = X \setminus A \in \mathcal{F}$
          (closed under complements);
    \item if $\{A_n\}_{n=1}^\infty \subseteq \mathcal{F}$, then
          $\bigcup_{n=1}^\infty A_n \in \mathcal{F}$ (closed under countable
          unions).
\end{enumerate}
The elements of $\mathcal{F}$ are called \emph{measurable sets}.
\end{definition}

\begin{definition}[Measurable space]
\label{def:app:measurable-space}
A \emph{measurable space} is a pair $(X, \mathcal{F})$ where $X$ is a nonempty
set and $\mathcal{F}$ is a $\sigma$-algebra on $X$.
\end{definition}

\begin{definition}[Borel $\sigma$-algebra]
\label{def:app:borel-sigma-algebra}
Let $(X, \tau)$ be a topological space. The \emph{Borel $\sigma$-algebra}
$\mathcal{B}(X)$ is the smallest $\sigma$-algebra containing all open sets of $X$.
Its elements are called \emph{Borel sets}.
\end{definition}

In particular, for $\R^n$ the Borel $\sigma$-algebra $\mathcal{B}(\R^n)$ is
generated by the open rectangles, or equivalently by the half-spaces of the form
$\{x \in \R^n : w^T x + b > 0\}$.

\begin{definition}[Measure]
\label{def:app:measure}
Let $(X, \mathcal{F})$ be a measurable space. A \emph{measure} is a function
$\mu \colon \mathcal{F} \to [0, +\infty]$ such that:
\begin{enumerate}
    \item $\mu(\emptyset) = 0$;
    \item \textbf{$\sigma$-additivity}: for any sequence
          $\{A_n\}_{n=1}^\infty$ of pairwise disjoint sets in $\mathcal{F}$,
          \[
          \mu\!\left(\bigcup_{n=1}^\infty A_n\right) = \sum_{n=1}^\infty \mu(A_n) .
          \]
\end{enumerate}
\end{definition}

\begin{definition}[Measure space]
\label{def:app:measure-space}
A \emph{measure space} is a triple $(X, \mathcal{F}, \mu)$ where $X$ is a nonempty
set, $\mathcal{F}$ is a $\sigma$-algebra on $X$, and
$\mu \colon \mathcal{F} \to [0, +\infty]$ is a measure. If $\mu(X) < \infty$, we
say $\mu$ is a \emph{finite measure}. If $\mu(X) = 1$, we say $\mu$ is a
\emph{probability measure}.
\end{definition}

\begin{definition}[Signed finite measure]
\label{def:app:signed-measure}
Let $(X, \mathcal{F})$ be a measurable space. A \emph{signed finite measure} is a
function $\mu \colon \mathcal{F} \to \R$ with $\mu(\emptyset) = 0$ that is
$\sigma$-additive: for any sequence $\{A_n\}_{n=1}^\infty$ of pairwise disjoint
sets in $\mathcal{F}$,
\[
\mu\!\left(\bigcup_{n=1}^\infty A_n\right) = \sum_{n=1}^\infty \mu(A_n) ,
\]
where the series converges absolutely. The \emph{total variation} of $\mu$ is the
measure $\abs{\mu}$ defined by
\[
\abs{\mu}(A) = \sup\!\left\{ \sum_{i=1}^n \abs{\mu(A_i)} :
\{A_i\}_{i=1}^n \text{ a finite partition of } A \right\} .
\]
\end{definition}

\begin{definition}[Lebesgue measure]
\label{def:app:lebesgue-measure}
The \emph{Lebesgue measure} $\lambda$ on $\R^n$ is the unique measure on
$(\R^n, \mathcal{B}(\R^n))$ that extends the notion of $n$-dimensional volume and
satisfies:
\begin{itemize}
    \item for a rectangle $R = [a_1, b_1] \times \cdots \times [a_n, b_n]$,
          $\lambda(R) = \prod_{i=1}^n (b_i - a_i)$;
    \item it is translation invariant: $\lambda(A + x) = \lambda(A)$ for all
          $x \in \R^n$.
\end{itemize}
\end{definition}

A set $A$ has \emph{Lebesgue measure zero} (or measure zero) if $\lambda(A) = 0$.
Examples include isolated points, finite sets, countable sets such as $\Q$, and
the Cantor set.

\begin{definition}[Measurable function]
\label{def:app:measurable-function}
Let $(X, \mathcal{F})$ and $(Y, \mathcal{G})$ be measurable spaces. A function
$f \colon X \to Y$ is \emph{measurable} (or $\mathcal{F}$-$\mathcal{G}$-measurable)
if the preimage of every measurable set is measurable:
\[
f^{-1}(B) = \{x \in X : f(x) \in B\} \in \mathcal{F}
\quad \forall B \in \mathcal{G} .
\]
\end{definition}

When $Y = \R$ with the Borel $\sigma$-algebra, it suffices to verify that
$\{x : f(x) > \alpha\} \in \mathcal{F}$ for every $\alpha \in \R$. In particular,
every continuous function is measurable with respect to the Borel
$\sigma$-algebras.

\begin{definition}[Lebesgue integral]
\label{def:app:lebesgue-integral}
Let $(X, \mathcal{F}, \mu)$ be a measure space. The \emph{Lebesgue integral} is
constructed in three steps:
\begin{enumerate}
    \item \textbf{Simple functions}: for $s = \sum_{i=1}^n a_i \chi_{A_i}$ with
          $a_i \geq 0$ and $A_i \in \mathcal{F}$ disjoint, where $\chi_A$ is the
          characteristic function of $A$, set
          \[
          \int_X s \, d\mu = \sum_{i=1}^n a_i \mu(A_i) .
          \]
    \item \textbf{Nonnegative functions}: for $f \colon X \to [0, +\infty]$
          measurable,
          \[
          \int_X f \, d\mu = \sup\!\left\{\int_X s \, d\mu :
          s \text{ simple, } 0 \leq s \leq f\right\} .
          \]
    \item \textbf{General functions}: for $f \colon X \to \R$ measurable, write
          $f = f^+ - f^-$ with $f^+ = \max(f, 0)$ and $f^- = \max(-f, 0)$. We say
          $f$ is \emph{integrable} (or $f \in L^1(\mu)$) if
          $\int_X f^+ \, d\mu < \infty$ and $\int_X f^- \, d\mu < \infty$, and set
          \[
          \int_X f \, d\mu = \int_X f^+ \, d\mu - \int_X f^- \, d\mu .
          \]
\end{enumerate}
\end{definition}

\begin{definition}[Almost everywhere]
\label{def:app:almost-everywhere}
Let $(X, \mathcal{F}, \mu)$ be a measure space. A property $P(x)$ holds
\emph{almost everywhere} (a.e.) if the set of points where it fails has measure
zero:
\[
\mu(\{x \in X : \neg P(x)\}) = 0 .
\]
\end{definition}

In the context of Lebesgue measure on $\R^n$, ``almost everywhere'' means
``except on a set of Lebesgue measure zero''.

\subsection{Functional-analytic preliminaries}
\label{subsec:app:functional-analysis}

\begin{definition}[Vector space]
\label{def:app:vector-space}
A \emph{vector space} over $\R$ is a set $V$ equipped with two operations,
addition $+ \colon V \times V \to V$ and scalar multiplication
$\cdot \colon \R \times V \to V$, satisfying the following axioms for all
$u, v, w \in V$ and $\alpha, \beta \in \R$:
\begin{enumerate}
    \item \textbf{Associativity of addition:} $(u + v) + w = u + (v + w)$;
    \item \textbf{Commutativity of addition:} $u + v = v + u$;
    \item \textbf{Identity element:} there exists $0 \in V$ with $v + 0 = v$;
    \item \textbf{Inverse element:} for each $v$ there is $-v \in V$ with
          $v + (-v) = 0$;
    \item \textbf{Compatibility of scalar multiplication:}
          $\alpha(\beta v) = (\alpha\beta)v$;
    \item \textbf{Unit element:} $1 \cdot v = v$;
    \item \textbf{Distributivity over vector addition:}
          $\alpha(u + v) = \alpha u + \alpha v$;
    \item \textbf{Distributivity over scalar addition:}
          $(\alpha + \beta)v = \alpha v + \beta v$.
\end{enumerate}
\end{definition}

\begin{definition}[Normed vector space]
\label{def:app:normed-space}
A \emph{normed vector space} is a pair $(X, \norm{\cdot})$ where $X$ is a real
vector space and $\norm{\cdot} \colon X \to [0, +\infty)$ is a \emph{norm}, that
is, a function satisfying:
\begin{enumerate}
    \item $\norm{x} = 0$ if and only if $x = 0$;
    \item $\norm{\alpha x} = \abs{\alpha} \norm{x}$ for all $\alpha \in \R$;
    \item $\norm{x + y} \leq \norm{x} + \norm{y}$ (triangle inequality).
\end{enumerate}
\end{definition}

Fundamental examples include $\R^n$ with the Euclidean norm and the function
spaces $L^p(\Omega)$.

\begin{definition}[Cover]
\label{def:app:cover}
Let $X$ be a topological space and $K \subseteq X$. A \emph{cover} of $K$ is a
family of sets $\{U_\alpha\}_{\alpha \in A}$ with
$K \subseteq \bigcup_{\alpha \in A} U_\alpha$. If each $U_\alpha$ is open, it is
an \emph{open cover}. A \emph{subcover} is a subfamily that is still a cover of
$K$.
\end{definition}

\begin{definition}[Compact space]
\label{def:app:compact-space}
A subset $K$ of a topological space is \emph{compact} if every open cover of $K$
admits a finite subcover.
\end{definition}

Intuitively, compactness generalizes the properties of closed and bounded sets in
$\R^n$ to general topological spaces.

\begin{theorem}[Heine--Borel]
\label{thm:app:heine-borel}
A subset $K \subset \R^n$ is compact if and only if it is closed and bounded
\citep[Theorem~2.41]{Rudin1976}.
\end{theorem}

In particular, the unit hypercube $I_m = [0,1]^m$ is compact, which is essential
for the universal approximation theorem.

\begin{definition}[Linear functional]
\label{def:app:linear-functional}
Let $X$ be a real vector space. A \emph{linear functional} is a map
$f \colon X \to \R$ satisfying
\[
f(\alpha x + \beta y) = \alpha f(x) + \beta f(y)
\quad \forall x, y \in X, \, \forall \alpha, \beta \in \R .
\]
That is, a linear functional is an element of the algebraic dual space $X^*$.
\end{definition}

\begin{definition}[Continuous linear functional]
\label{def:app:continuous-functional}
Let $(X, \norm{\cdot})$ be a normed vector space. A linear functional
$f \colon X \to \R$ is \emph{continuous} (or \emph{bounded}) if there is a
constant $M > 0$ with
\[
\abs{f(x)} \leq M\norm{x} \quad \forall x \in X .
\]
\end{definition}

The set of all continuous linear functionals on $X$ forms the \emph{dual space}
$X'$, a Banach space under the operator norm
$\norm{f} = \sup_{\norm{x} \leq 1} \abs{f(x)}$.

\begin{definition}[The space $C(K)$]
\label{def:app:space-CK}
Let $K$ be a compact space. The space $C(K)$ is the set of all continuous
functions $f \colon K \to \R$ equipped with the \emph{supremum norm} (or uniform
norm)
\[
\norm{f}_\infty = \sup_{x \in K} \abs{f(x)} .
\]
\end{definition}

With this norm, $C(K)$ is a Banach space (a complete normed vector space). In the
context of the universal approximation theorem, we work with $C(I_m)$, the space
of continuous functions on the unit hypercube.

\begin{theorem}[Lebesgue dominated convergence]
\label{thm:app:dominated-convergence}
Let $(X, \mathcal{F}, \mu)$ be a measure space and $\{f_n\}_{n=1}^\infty$ a
sequence of measurable functions $f_n \colon X \to \R$ such that:
\begin{enumerate}
    \item $f_n(x) \to f(x)$ for almost every $x \in X$;
    \item there is an integrable function $g \colon X \to \R$ with
          $\abs{f_n(x)} \leq g(x)$ for almost every $x \in X$ and every
          $n \in \mathbb{N}$.
\end{enumerate}
Then $f$ is integrable and
\[
\lim_{n \to \infty} \int_X f_n \, d\mu = \int_X f \, d\mu
= \int_X \lim_{n \to \infty} f_n \, d\mu
\]
\citep[Theorem~2.24]{Folland1999}.
\end{theorem}

This theorem is fundamental for interchanging limits and integrals, which is
needed to show that sigmoidal functions are discriminatory.

\begin{definition}[Convex set]
\label{def:app:convex-set}
A subset $C$ of a vector space $V$ is \emph{convex} if for all $x, y \in C$ and
$\lambda \in [0,1]$ one has $\lambda x + (1-\lambda)y \in C$. Equivalently, $C$
contains the segment joining any two of its points.
\end{definition}

\begin{theorem}[Hahn--Banach, geometric form]
\label{thm:app:hahn-banach}
Let $X$ be a real topological vector space, and $A$, $B$ disjoint convex subsets
of $X$ with $A$ open. Then there exist a continuous linear functional
$f \colon X \to \R$ and a constant $\alpha \in \R$ such that
\[
f(a) < \alpha \leq f(b) \quad \forall a \in A, \, b \in B
\]
\citep[Theorem~5.5]{Folland1999}.
\end{theorem}

\begin{corollary}[Density characterization]
\label{cor:app:density-characterization}
Let $X$ be a normed vector space and $Y \subset X$ a subspace. Then $Y$ is dense
in $X$ if and only if the only continuous linear functional vanishing on $Y$ is
the zero functional.
\end{corollary}

\begin{proof}
\textbf{($\Rightarrow$)} Suppose $Y$ is dense in $X$. Let $f \colon X \to \R$ be
a continuous linear functional with $f(y) = 0$ for all $y \in Y$. For any
$x \in X$ there is a sequence $\{y_n\} \subset Y$ with $y_n \to x$. By continuity
of $f$,
\[
f(x) = f\!\left(\lim_{n \to \infty} y_n\right)
= \lim_{n \to \infty} f(y_n) = 0 ,
\]
so $f \equiv 0$.

\textbf{($\Leftarrow$)} Suppose $Y$ is not dense in $X$. Then $\overline{Y} \neq X$,
and there is $x_0 \in X \setminus \overline{Y}$. Since $\overline{Y}$ is a closed
subspace, take an open ball $B_r(0) \subset X$ and set
$A = \overline{Y} + B_r(0)$, which is open and convex with $A \cap \{x_0\} =
\emptyset$ for $r$ small enough (as $x_0 \notin \overline{Y}$). By
\Cref{thm:app:hahn-banach}, there exist a linear functional $f$ and
$\alpha \in \R$ with
\[
f(a) < \alpha \leq f(x_0) \quad \forall a \in A .
\]
Since $\overline{Y}$ is a subspace, $0 \in \overline{Y} \subset A$, so
$f(0) = 0 < \alpha$. Moreover, for any $y \in Y$ and $t \in \R$, $ty \in
\overline{Y} \subset A$, so $tf(y) = f(ty) < \alpha$ for all $t$. This is
possible only if $f(y) = 0$. Yet $f(x_0) \geq \alpha > 0$, so $f \not\equiv 0$.
\end{proof}

\begin{theorem}[Riesz representation for $C(K)$]
\label{thm:app:riesz}
Let $K$ be a compact space and $C(K)$ the space of real continuous functions on
$K$ with the supremum norm. For every continuous linear functional
$L \colon C(K) \to \R$ there exists a unique regular finite signed Borel measure
$\mu$ on $K$ such that
\[
L(f) = \int_K f \, d\mu \quad \forall f \in C(K) .
\]
Moreover, $\norm{L} = \abs{\mu}(K)$, where $\abs{\mu}$ is the total variation of
$\mu$ \citep[Theorem~7.17]{Folland1999}.
\end{theorem}

\begin{theorem}[Stone--Weierstrass]
\label{thm:app:stone-weierstrass}
Let $K$ be a compact space and $\mathcal{A} \subset C(K)$ a subalgebra that
separates points of $K$ and vanishes nowhere. Then $\mathcal{A}$ is dense in
$C(K)$ with respect to the supremum norm \citep[Theorem~7.32]{Rudin1976}.
\end{theorem}

In particular, if $\mathcal{A}$ contains the constant functions and separates
points, then it is dense in $C(K)$.

\begin{remark}
There are two main approaches to proving the universal approximation theorem. The
approach of \citet{cybenko1989} uses the Hahn--Banach corollary together with the
Riesz representation theorem, showing that sigmoidal functions are
``discriminatory'' (they vanish against a measure only if the measure is zero).
The approach of \citet{hornik1989} uses the Stone--Weierstrass theorem, showing
that neural networks form a point-separating subalgebra. Both rely on the
preliminaries collected above, and the full statement and proof are developed in
\Cref{ch:foundations}.
\end{remark}
